\chapter[State of the Art]{State of the Art and Technical Background}

Radio map prediction sits at the intersection of classical wireless propagation
modelling, statistical channel analysis, and dense computer vision regression.
This chapter reviews the background needed to understand the final design:
classical and empirical propagation models, deep learning for radio map
estimation, physics informed and hybrid neural approaches, distribution aware
prediction, and fair comparison rules. The goal is to identify the limitations
of existing work and the tools that will be used later, not to report the
results of this thesis.

% Section separator
\section{General Framework: CKM Surrogates and Evaluation Contracts}
\label{sec:soa_ckm_context}
% Section separator

The term channel knowledge map (CKM) refers to a site specific database of
location tagged channel knowledge, such as channel gain, angles, or delay
information, used to make future radio systems environment aware
\cite{zengx2021ckm,ckmtutorial2024}. This framing is more precise than a
generic radio map: the map is not only an image like coverage estimate, but a
queryable prior that can reduce channel state acquisition overhead, guide
beamforming, support planning, and provide initial channel knowledge before
online measurements are available.

Recent CKM literature also clarifies why the thesis uses a large offline
ray traced dataset rather than sparse field samples. Xu and Zeng
\cite{ckmdata2024} analyse how CKM construction error depends on spatial sample
density and on the number of neighbouring samples used for online prediction,
showing the explicit cost/accuracy trade off of measurement driven CKM
construction. For a thesis whose goal is to learn dense \(513\times513\)
targets under city holdout, this motivates the simulation first route: it
provides full supervision everywhere, while any learned model is still best
understood as a surrogate for a ray traced CKM rather than as a replacement for
field calibration.

The data ecosystem has also expanded since the first path loss challenge.
RadioMap3DSeer \cite{dataset2212} and the ICASSP~2023 challenge
\cite{icassp2023challenge} provide the closest open benchmark family for
path loss image prediction, while the directive antenna dataset of Jaensch et
al.\ \cite{jaensch2024directiverme} and CKMImageNet
\cite{ckmimagenet2025} show the broader move toward open, location specific
datasets with richer channel and environment representations. The CKM dataset
used in this thesis sits in the same data driven tradition, but differs in its
continuous UAV height conditioning, real city holdout split, and simultaneous
channel attenuation, delay spread, and angular spread targets.

% Section separator
\section{Path Loss Prediction}
\label{sec:soa_classical}
% Section separator

\subsection{Free space and two ray propagation}

The most fundamental path loss term is the \emph{free space path loss} (FSPL),
which depends only on carrier frequency~$f$ and Tx/Rx distance~$d$:
\begin{equation}
  \mathrm{PL}_{\mathrm{FSPL}} = 20\log_{10}(d) + 20\log_{10}(f)
    + 20\log_{10}\!\left(\frac{4\pi}{c}\right).
\end{equation}
FSPL is the single direct ray reference baseline in unobstructed conditions, not
a lower bound on the final path loss once coherent reflections are present. In
aerial scenarios a specular ground reflection can produce constructive or
destructive interference periodically with distance. Let \(d_{2D}\) be the
horizontal Tx/Rx distance, \(h_{\mathrm{tx}}\) the UAV height, and
\(h_{\mathrm{rx}}\) the receiver height. The direct and reflected path lengths
are
\begin{equation}
  d_{\mathrm{dir}} =
  \sqrt{d_{2D}^{2} + (h_{\mathrm{tx}}-h_{\mathrm{rx}})^{2}},
  \qquad
  d_{\mathrm{ref}} =
  \sqrt{d_{2D}^{2} + (h_{\mathrm{tx}}+h_{\mathrm{rx}})^{2}} .
\end{equation}
With wavelength \(\lambda\), the coherent two ray model adds the direct and
ground reflected fields before converting to dB \cite{rappaport,wocc2021}:
\begin{equation}
  E_{\mathrm{total}}
  =
  \frac{e^{-j2\pi d_{\mathrm{dir}}/\lambda}}{d_{\mathrm{dir}}}
  +
  \Gamma
  \frac{e^{-j2\pi d_{\mathrm{ref}}/\lambda}}{d_{\mathrm{ref}}},
\end{equation}
where \(\Gamma\) is the effective ground reflection coefficient. The received
field magnitude \(|E_{\mathrm{total}}|\) is then converted to a path loss like
dB prediction with a fitted calibration constant. This is the physical origin
of the radial rings later exploited by the path loss prior.

\begin{figure}[ht]
  \centering
  \begin{tikzpicture}[x=0.92cm,y=0.92cm,>=Latex,font=\small]
    \coordinate (Tx) at (0.2,3.25);
    \coordinate (Rx) at (6.6,0.45);
    \coordinate (Gtx) at (0.2,0);
    \coordinate (Grx) at (6.6,0);
    \coordinate (RxImg) at (6.6,-0.45);
    \coordinate (R) at (5.82,0);

    \draw[thick] (-0.75,0) -- (7.35,0) node[below right] {ground};
    \draw[densely dotted] (Tx) -- (Gtx);
    \draw[densely dotted] (Rx) -- (RxImg);

    \draw[<->] (-0.35,0) -- node[left] {\(h_{\mathrm{tx}}\)} (-0.35,3.25);
    \draw[<->] (7.02,0) -- node[right] {\scriptsize \(h_{\mathrm{rx}}\)} (7.02,0.45);
    \draw[<->] (0.2,-0.78) -- node[below] {\(d_{2D}\)} (6.6,-0.78);

    \filldraw[blue] (Tx) circle (2pt) node[above left] {UAV Tx};
    \filldraw[black] (Rx) circle (2pt) node[above right] {ground Rx};
    \filldraw[gray] (RxImg) circle (1.3pt)
      node[below right,gray] {\scriptsize image Rx};
    \filldraw[gray] (R) circle (1.4pt);

    \draw[->,thick] (Tx) --
      node[above,pos=0.60,sloped] {direct path} (Rx);
    \draw[densely dotted,gray] (Tx) -- (RxImg);
    \draw[->,thick,dashed] (Tx) -- (R);
    \draw[->,thick,dashed] (R) -- (Rx);
    \node[fill=white,inner sep=1pt] at (3.1,1.25) {ground reflected path};
    \node[below,fill=white,inner sep=1pt] at (R) {\scriptsize reflection point};
  \end{tikzpicture}
  \caption{Direct and specular ground reflected paths in the two ray model.}
  \label{fig:two_ray_schematic}
\end{figure}

\subsection{Empirical and semi empirical models}

Empirical models such as the Okumura-Hata and COST-231/Hata formulas extend
FSPL by incorporating frequency, antenna heights, and city type correction terms
\cite{rappaport,cost231}.  While computationally cheap, they are calibrated on specific
frequency bands and city/environment types, limiting accuracy when the geometry
departs from the fitting dataset.

The 3GPP Technical Report 38.901 \cite{tr38901} formalises stochastic channel
models for frequencies from 500~MHz to 100~GHz for urban macro (UMa), urban
micro (UMi), rural macro (RMa), and indoor scenarios.  It provides separate
path loss expressions for LoS and non-LoS (NLoS) conditions, breakpoint
distances, and shadow fading standard deviations, which serve as reference
baselines throughout this thesis.

\subsection{UAV and air to ground channel models}

Classical terrestrial channel models assume the transmitter is at rooftop height
or below.  Unmanned aerial vehicle (UAV) base station scenarios introduce a
third dimension: the UAV can fly from near ground level to several hundred
metres.  Al-Hourani et al.~\cite{alhourani2014} showed that the LoS probability
increases monotonically with elevation angle and depends on city/environment type
(suburban, urban, dense urban, high rise).  The A2G modelling literature
\cite{alhourani2014,khawaja_survey} commonly parameterises this dependence
through logistic functions of the elevation angle
$\theta = \arctan\!\bigl((h_{\mathrm{tx}} - h_{\mathrm{rx}})/d_{2D}\bigr)$,
where $h_{\mathrm{tx}}$ and $h_{\mathrm{rx}}$ are transmit and receive heights
respectively and $d_{2D}$ is the horizontal distance. 3GPP~TR~38.901
\cite{tr38901} is used here as standardized LoS/NLoS large scale channel model
context, not as the direct source of that elevation-logistic A2G formula.

Recent work by Saboor and Vinogradov \cite{saboor2025height} makes this
height dependence especially relevant for the present thesis: in simulated
26~GHz urban A2G links, LoS path loss remains close to a free space exponent,
whereas NLoS path loss exponents and shadow fading evolve with ABS height and
with urban layout. Vinogradov et al.\ \cite{vinogradov2026shadow} also propose
3D shadow projections as a fast way to generate spatially consistent A2G LoS
maps without full point by point ray tracing. These results support two design
choices used throughout this work: treating the LoS/NLoS mask as a physical
input rather than as a cosmetic label, and conditioning every dense prediction
on UAV height.

These works motivate a useful separation for the thesis methodology: LoS/NLoS
visibility and continuous UAV height should be treated as physical conditioning
variables, not as decorative labels appended to a generic image regression
network.

\subsection{Spread definitions used later}

Beyond path loss, full channel characterisation requires second order statistics.
The \emph{root mean square delay spread} $\tau_{\mathrm{rms}}$ measures the
time domain dispersion of multipath arrivals and drives inter symbol interference
in wideband systems.  The \emph{azimuth angular spread} $\sigma_{\mathrm{AS}}$
at the receiver quantifies the spatial diversity of arriving paths.

For multipath components with powers \(P_\ell\), delays \(\tau_\ell\), and
azimuth angles \(\phi_\ell\), the RMS delay spread is
\begin{equation}
  \bar{\tau} =
  \frac{\sum_\ell P_\ell \tau_\ell}{\sum_\ell P_\ell},
  \qquad
  \tau_{\mathrm{rms}} =
  \sqrt{
  \frac{\sum_\ell P_\ell(\tau_\ell-\bar{\tau})^2}{\sum_\ell P_\ell}
  } .
\end{equation}
For angular spread, the mean angle must be computed circularly:
\begin{equation}
  \bar{\phi} =
  \operatorname{atan2}\!\left(
    \sum_\ell P_\ell \sin\phi_\ell,
    \sum_\ell P_\ell \cos\phi_\ell
  \right),
  \qquad
  \sigma_{\mathrm{AS}} =
  \sqrt{
  \frac{\sum_\ell P_\ell\,\operatorname{wrap}(\phi_\ell-\bar{\phi})^2}
       {\sum_\ell P_\ell}
  } ,
\end{equation}
where \(\operatorname{wrap}(\cdot)\) maps angle differences to the
\((-\pi,\pi]\) interval.

3GPP~TR~38.901 \cite{tr38901} models both quantities as log normal random
variables conditioned on the LoS/NLoS state and distance.  WINNER~II follows
the same broad log domain large scale parameter convention \cite{winner2}.
In practice, the distributions of delay spread and angular spread are
right skewed with a heavy exponential tail and often include a spike near zero
representing the dominant LoS path. A Gaussian model cannot represent this
shape faithfully. This is why the methodology later checks the empirical target
histograms before choosing distribution aware heads.

% Section separator
\subsection{Deep learning for path loss maps}
\label{sec:soa_dl}
% Section separator

\subsubsection{Image to image translation baselines}

The radio map prediction problem can be cast as an image to image regression:
given an input topology map (and optionally a transmitter location channel), the
model predicts a path loss map of the same spatial extent.  The pix2pix framework
\cite{isola2017pix2pix} established a general recipe for such tasks using a
conditional GAN with a U-Net generator and a PatchGAN discriminator.  Early
iterations of this thesis used this formulation.  In practice the adversarial
objective sharpens spatial texture but does not address the height dependent
radial pattern or the heavy NLoS tails that dominate the error budget in the CKM
dataset.

RadioUNet \cite{radiounet2020} is the canonical convolutional path loss map
baseline before the PMNet challenge family. In its accurate map settings it
reports grey level RMSE values around \(0.0203\)--\(0.0384\); using the paper's
\SI{80}{dB} grey level conversion gives an approximate \SI{1.6}{dB}--\SI{3.1}{dB}
path loss scale. RadioGUNet \cite{radiogunet2025} extends this outdoor
RadioMapSeer line with group equivariant convolutions, reporting
\SI{1.304}{dB} on the easiest DPM setting, \SI{1.936}{dB} on the IRT setting,
and \SI{1.392}{dB} with cars. These values are used again in
Chapter~\ref{chap:results} as scale references for the final channel attenuation result.

\textbf{Limitations:} These papers are useful path loss scale references, but
they do not report dense delay spread or angular spread maps and do not use the
CKM city holdout, continuous UAV height contract.

\subsubsection{PMNet and the ICASSP outdoor challenge family}

PMNet \cite{pmnet_icassp2023} was the winner of the first ICASSP 2023 Outdoor Pathloss
Radio Map Prediction Challenge \cite{icassp2023challenge}.  It uses a deep
encoder decoder with skip connections, trained and evaluated on the
RadioMap3DSeer dataset~\cite{dataset2212}.  RadioMap3DSeer consists of
city maps derived from urban 3D geometry with rooftop transmitters placed
3~m above the selected rooftop in the dataset description, and building
interior pixels masked to zero error.

The challenge metric is RMSE in normalised image units.  As an approximate
dB scale reference, the RadioMap3DSeer image mapping clips path gain to the range
$[\mathrm{PL}_{\mathrm{thr}}, \mathrm{PL}_{\mathrm{max}}] = [-111, -75]\,\mathrm{dB}$,
yielding a $36\,\mathrm{dB}$ dynamic window under that interpretation.
Multiplying normalised RMSE by~36 gives the rough dB scale references listed in
Table~\ref{tab:icassp2023corrected}.

\begin{table}[ht]
\centering
\caption[Approximate ICASSP 2023 outdoor RMSE scale.]%
{Approximate dB scale RMSE references for ICASSP~2023 outdoor challenge methods.
The normalised scores are taken from the challenge family and converted with the
\SI{36}{dB} RadioMap3DSeer image window described in the dataset
paper \cite{icassp2023challenge,dataset2212}.\protect\footnotemark}
\label{tab:icassp2023corrected}
\begin{tabular}{lcc}
\toprule
\textbf{Method} & \textbf{RMSE (norm.)} & \textbf{Approx. RMSE (dB)} \\
\midrule
PMNet \cite{pmnet_icassp2023} & 0.0383 & 1.38 \\
PPNet baseline \cite{icassp2023challenge} & 0.0507 & 1.83 \\
\bottomrule
\end{tabular}
\end{table}
\footnotetext{This conversion should be read as an approximate scale reference,
not as a like for like physical RMSE. It uses the path gain
clipping window $[\mathrm{PL}_{\mathrm{thr}},\mathrm{PL}_{\mathrm{max}}] =
[-111,-75]\,\mathrm{dB}$ associated with the normalised target image, so the
physical span is $-75-(-111)=36\,\mathrm{dB}$. Multiplying the normalised RMSE
values reported in the challenge paper \cite{icassp2023challenge,dataset2212}
by this factor converts them to physical dB, which gives
$0.0383\times 36\approx 1.38$ for PMNet.}

\textbf{Limitations:} These values are scale references only. They are not
directly comparable to CKM because RadioMap3DSeer uses a different dataset,
fixed transmitter placement, a much narrower target image window, and a
different building pixel convention. The CKM task later evaluated in this
thesis uses unseen real city geometries, native \(513\times513\) maps,
continuous UAV altitude, and metrics over valid ground receiver pixels. The
challenge includes multiple transmitter locations, but it does not test the
continuous UAV height conditioning used in CKM.

\subsubsection{Transformers, state space models, and attention based encoders}

The RMTransformer \cite{rmtransformer2025} evaluates a hybrid MaxViT/CNN
encoder--decoder on the USC/PMNet dataset rather than on the ICASSP challenge
test set.  It reports a normalised RMSE of~0.007148 versus~0.01046 for PMNet
under a random 90/10 split on \(256\times256\) maps.  The paper normalises
received power/path loss from \(-254\) to \SI{0}{dBm}, so the reported pixel
RMSE corresponds to roughly \SI{1.82}{dB} on its own scale, and the reported
channel prediction error of 0.008099 corresponds to roughly \SI{2.06}{dB}.
This supports the use of attention for long range urban context, but it is not
inserted in Table~\ref{tab:icassp2023corrected} because it is not measured on
RadioMap3DSeer or under the ICASSP challenge contract.

Recent attention and foundation model variants extend the effective receptive
field of radio map predictors \cite{wicopg2025,fmrme2026}. This is useful when
large scale urban context affects the prediction, but these methods normally
need more data and compute than convolutional models and their reported
benchmarks do not match the CKM UAV height contract.

\subsubsection{Foundation models for radio maps}

The most recent shift in the literature is towards \emph{wireless channel
foundation models} pretrained on large heterogeneous datasets before
task specific fine tuning.  WiCo-PG \cite{wicopg2025} uses dual VQGANs and a
transformer with frequency-guided mixture-of-experts.  It exploits RGB drone
camera images as an auxiliary modality and achieves an NMSE of~0.012 on its
constructed multimodal U2G benchmark.  RadioLAM \cite{radiolam2025} targets
fine grained 3D radio map estimation from ultra low sampling rates using a
large generative model pipeline.  FM-RME \cite{fmrme2026} frames radio map estimation as a foundation
modelling problem with self supervised pretraining on simulated propagation
data.

\textbf{Limitations:} These foundation approaches are relevant as long term
background, but they require training corpora and compute budgets beyond a
single CKM dataset. They are therefore cited as future context, not as methods
directly compared with the final thesis model.

% Section separator
\subsection{Physics informed and hybrid path loss methods}
\label{sec:soa_physics}
% Section separator

\subsubsection{Geometry assisted and diffraction aware deep learning}

Purely data driven models must implicitly learn diffraction geometry from the
topology map.  An alternative is to precompute a physics based channel prior and
provide it as an additional input channel, reducing the learning problem to
residual estimation.

Shadow projection methods \cite{vinogradov2026shadow} occupy an intermediate
point between stochastic LoS probability models and full ray tracing. They use
3D building geometry to project blocked ground regions from the aerial
transmitter, producing a deterministic LoS map that is spatially consistent
along user trajectories. The CKM dataset already provides a ray traced LoS
mask, so this thesis does not need to generate the mask itself; nevertheless,
the method is directly relevant because it explains why a binary visibility
channel can carry so much of the structure that the neural network otherwise
has to infer from the height map alone.

Geometry assisted diffraction features~\cite{geomDL2024} use virtual obstacles,
multi screen knife edge propagation, and scattering aware local geometry cues to
inject physically meaningful NLoS information into a deep model.  This is
closely related to the motivation behind the diffraction-prior experiments in
this thesis, although the implementation used later is a lightweight per map
channel rather than the full geometry model assisted learning framework of that
paper. The limitation is that a single precomputed diffraction cue cannot
replace a calibrated NLoS model when dense urban shadowing dominates.

\subsubsection{Physics informed neural networks}

The ReVeal architecture \cite{reveal2025} operates as a physics informed neural
network (PINN) by incorporating a second order PDE residual as an auxiliary loss
term.  This penalises predictions with inconsistent local spatial curvature
under the derived propagation constraint, reducing corner transition artefacts.
ReVeal achieves~$1.95\,\mathrm{dB}$ RMSE in rural scenarios using only 30
training points by leveraging the physical constraint as a strong regulariser.
A lightweight analogue for dense radio maps is a masked Laplacian penalty on
the residual prediction, which penalises locally nonsmooth outputs without
solving the full ReVeal operator.

\subsubsection{Prior plus residual correction}

A principled hybrid approach is to \emph{freeze} a calibrated analytic prior and
train the neural component to predict only the structured residual.  This design
reduces the effective learning problem from the full path loss range to a
narrower residual distribution, improving data efficiency and interpretability.

Chapter~3 applies this philosophy to the three thesis targets. For the channel
attenuation target, an enhanced two ray path loss model fitted per height bin provides a prior map
$\hat{y}_{\mathrm{PL}}(\textbf{x}) = \mathrm{TwoRay}(d, h_{\mathrm{tx}})$,
and the neural component predicts the residual
$\epsilon = y - \hat{y}_{\mathrm{PL}}$. For delay spread and angular spread,
the same prior residual idea can be implemented with a ridge regression model
in log space using TX centred geometric features and local morphology
statistics. The important background point is the decomposition itself: a
frozen train calibrated prior defines what the neural model should not have to
rediscover.

% Section separator
\section{Angular Spread Prediction}
\label{sec:soa_stats}
% Section separator

\subsection{Empirical models for spread parameters}

Standard channel models characterise delay spread and angular spread through
log domain large scale parameters tied to propagation state and scenario.
TR~38.901 \cite{tr38901} provides scenario specific tables and formulae for
these parameters, with dependencies on variables such as distance, frequency,
terminal heights, and LoS/NLoS state across UMa, UMi, and RMa scenarios.  These
parameterisations are limited to the morphology class and frequency band of the
underlying measurement campaigns.

Log domain modelling is physically motivated: spread parameters are
multiplicative quantities in physical space, so their logarithm is more
Gaussian like and additive noise assumptions are better justified. The calibrated spread prior adopts
this convention by working in $\log(1+y)$ space throughout, consistent with the
3GPP empirical modelling framework.

\subsection{Deep learning for wideband channel statistics}

Work on deep learning for delay and angular spread is sparser than for path
loss. Most published methods either predict spread globally per sample (a single scalar)
rather than per pixel, or treat it as an ancillary output of a multitask
path loss model.  Per pixel dense prediction of delay spread and angular spread
over $513 \times 513$ maps is, to the best of the author's knowledge, not
addressed by published methods at the time of writing.

The closest external quantitative references for spreads are therefore scalar
channel characteristic predictors rather than dense CKM-map benchmarks. Huang et al.\
\cite{huang2025a2gtransformer} are not used in the results comparison tables
because their Wireless InSite campus dataset predicts scalar A2G link
characteristics from coordinates and link features, not dense CKM maps from
city topologies.

\textbf{Limitation for this thesis.}
Most reviewed angular spread work reports scalar link level channel parameters
or simulation specific statistics rather than dense CKM maps under city
holdout. Chapter~4 therefore compares angular spread mainly against the frozen
prior, the thesis target, and one or two papers with similar but not equal
goals.

The main challenge is the spike plus exponential distribution shape of spread
targets.  A model trained only with mean squared error (MSE) collapses to the
conditional mean, which lies in the interior of the spike and the exponential
tail and is therefore a poor representative of either regime. This motivates
Gaussian mixture or other distribution aware heads, especially when a
log domain prior is available as an anchor before residual learning.

\section{Delay Spread Prediction}
\label{sec:soa_delay_spread}

Delay spread has the same dense map gap as angular spread, but the physical
interpretation is temporal rather than directional. The target is small for
many LoS dominated receivers and can become large in streets or blocks with
multiple reflected and diffracted paths. This makes the quantity sensitive to
rare high delay pixels, so a mean-error metric can hide the tail behaviour that
matters for wideband links.

\textbf{Limitation for this thesis.}
Most reviewed delay spread work reports scalar link level channel parameters
or simulation specific statistics rather than dense \(513\times513\) CKM maps
under city holdout. Chapter~4 therefore compares delay spread mainly against
the frozen prior, the thesis target, and one or two papers with similar but not
equal goals.

\subsection{Distributional mismatch and the spike plus exponential problem}

The spike plus exponential shape, consisting of a Dirac like concentration near zero
(the LoS dominated case) plus a heavy exponential tail (NLoS multipath), was
confirmed empirically on the CKM dataset.  Skewness exceeds~5 and excess kurtosis
exceeds~40 for both delay and angular spread in several topology classes.  A
single Gaussian or even a symmetric loss function does not fit this family.

Mixture models are the natural solution, but the required number of components
depends on what is being modelled. A histogram study can select the best
family for the raw target distributions, where path loss and the spread targets
may need several components to match the full marginal shape. In a
prior residual setting, however, the mixture head models a different quantity:
the residual distribution after a frozen prior has already explained the
dominant LoS, NLoS, delay spread, and angular spread structure. This
distinction matters because the prior can remove much of the target side
multimodality, leaving a compact mixture to represent under correction,
near zero correction, and over correction modes.

% Section separator
\section{Training and Fair Comparison}
\label{sec:soa_distribution}
% Section separator

\subsection{Heteroscedastic uncertainty}

Kendall and Gal \cite{kendall2017uncertainties} established the framework for
heteroscedastic neural networks in computer vision: predict both a mean output
$\mu(x)$ and a log variance $\log \sigma^2(x)$, optimising the negative
log likelihood (NLL)
\begin{equation}
  \mathcal{L} = \frac{(\mu - y)^2}{\sigma^2} + \log \sigma^2.
\end{equation}
Pixels with high aleatoric uncertainty (deep NLoS shadows, building edges)
receive lower effective weight, while the network is penalised for
overestimating uncertainty everywhere. A PMHHNet uncertainty diagnostic applied this
loss to path loss prediction, with $\sigma^2$ predicted as a second output
channel of PMHHNet.

\subsection{Gaussian mixture model heads}

Beyond heteroscedastic Gaussians, the GMM head allows multimodal output
distributions.  Mixture density networks have already appeared in wireless
channel modelling: Saleh et al.\ \cite{saleh2021probabilistic} use an MDN to
predict probabilistic mmWave path loss distributions, Garcia Marti et al.\
\cite{garciamarti2020mixture} learn a stochastic Gaussian mixture channel
model for physical layer design, and Lee et al.\ \cite{lee2024timevarying}
derive conditional received power PDFs with a mixture density network.  These
works support the use of mixture outputs for wireless propagation, but they
operate at the level of scalar or channel sample distributions rather than
dense radio map generation.

A dense CKM version of a GMM head can be implemented with a two stage
architecture. Stage~A produces per sample GMM parameters
$\{\pi_k, \mu_k, \sigma_k\}_{k=1}^{K}$ via a convolutional encoder conditioned
on UAV height; Stage~B produces per pixel component assignments \(p_k(i,j)\)
and within component deviations \(z(i,j)\). The prediction is
\begin{equation}
  \hat{y}(i,j) = \sum_{k=1}^{K} p_k(i,j) \cdot \bigl(\mu_k + z(i,j) \cdot \sigma_k\bigr).
\end{equation}
This formulation avoids the mode collapse failure of MSE regression while
maintaining a compact architecture. The important idea is therefore not the
isolated use of a Gaussian mixture, but its role as a global distribution head
that tells the dense decoder what value distribution the map should contain
before the model solves the harder pixel placement problem.

\subsection{Stochastic diffusion models}

Conditional diffusion models applied to radio maps \cite{radiodiff2025} generate
physically plausible stochastic variations in deep shadow zones rather than
predicting a single blurry mean.  The iterative denoising procedure can be
conditioned on topology maps, transmitter locations, and obstacle prompts.
RadioDiff reports lower normalised RMSE and higher SSIM/PSNR than CNN and
sequence model baselines on static and dynamic RadioMapSeer settings, but its
metrics are image normalised rather than physical dB values.  These models are
computationally expensive at inference time and are listed as a future work
horizon for this thesis.

\subsection{Multitask and joint prediction}

Training a single shared backbone across multiple output targets can improve
sample efficiency and consistency between targets.  Channel attenuation, delay spread, and
angular spread are physically correlated: a deep shadow corresponds to high
delay spread and potentially wide angular spread from diffracted paths.  A joint
model can therefore use the intermediate feature maps learned for attenuation to
partially explain spread patterns, reducing the effective learning problem.

Chapter~3 uses this idea by sharing the topology and UAV height representation
across the three targets, while allowing target specific residual heads for
channel attenuation, delay spread, and angular spread. The literature point is that joint
training is useful only if the shared representation is anchored by
target specific priors and losses; otherwise the easiest target can dominate
the common backbone.

% Section separator
\subsection{Training stabilisation and regularisation techniques}
\label{sec:soa_training}
% Section separator

\noindent
The methods in this section are included because dense radio map prediction is
easy to overfit. The model sees complete city maps with strong spatial
correlations, while validation and test cities are completely held out. Small
training subsets, heavy tailed NLoS errors, and quantized targets can make the
training loss improve even when the deployment-relevant city holdout error does
not. Stabilisation and regularisation techniques are therefore practical
safeguards against selecting a model that only fits
the training cities.

\subsubsection{Stochastic weight averaging}

Stochastic weight averaging (SWA) \cite{izmailov2018swa} averages model weights
along the training trajectory at a low, approximately constant learning rate.
This finds wider optima in the loss landscape that generalise better to
unseen data. The CKM city holdout evaluation makes generalisation particularly
important since validation and test cities are never seen during training.
In this thesis, SWA was useful as a diagnostic stabilisation tool in earlier
neural experiments, although it is not part of every final training recipe.

\subsubsection{Exponential moving average weights}

Separate from SWA, exponential moving average (EMA) tracking of the model
weights at a decay rate of~0.9975 can produce smoother validation checkpoints.
Earlier diagnostic branches used EMA weights for validation scoring and fed a
smoothed score to the ReduceLROnPlateau scheduler. This is useful background
for interpreting the training choices in Chapter~3, where the final checkpoint
selection rule is stated explicitly for reproducibility.

\subsubsection{Feature wise linear modulation}

Feature wise linear modulation (FiLM) \cite{perez2018film} conditions
the activations of a neural network on a scalar or vector input by applying
learned affine transforms $\gamma(c) \cdot \textbf{h} + \beta(c)$ to
intermediate feature maps.  Sinusoidal positional encoding of the UAV height
scalar (using log spaced frequencies over the full 12 to 478~m range) is fed into
FiLM layers at multiple encoder depths.  This ensures that the UAV altitude
continuously modulates the activation statistics rather than being handled as a
discrete bin, which is critical when the effective propagation regime changes
substantially with height.

% Section separator
\subsection{Key benchmarks and fair comparison}
\label{sec:soa_benchmarks}
% Section separator

\subsubsection{Why headline numbers cannot be compared directly}

The most widespread mistake in radio map literature reviews is comparing
headline RMSE values without verifying whether the underlying tasks are
equivalent.  The following evaluation axes are the minimum required for a
valid comparison:

\begin{enumerate}
  \item \textbf{Split:} city holdout vs.\ in-distribution random split vs.\
        same-scene split.
  \item \textbf{Transmitter height:} fixed rooftop, discrete bins, or
        continuous UAV altitude.
  \item \textbf{Building pixels:} excluded from loss, excluded from metrics,
        set to zero error, or included.
  \item \textbf{Target quantity and units:} path loss (dB), path gain (dB),
        received power (dBm), delay spread (ns), or angular spread (degrees).
  \item \textbf{Dynamic range and clipping:} determines the RMSE floor and
        ceiling independently of model quality.
  \item \textbf{Inference extras and calibration source:} test time
        augmentation (TTA), SWA, EMA weights, and whether calibration is
        train only, validation based, or fitted from test/field measurements.
\end{enumerate}

This thesis uses a strict city holdout split, continuous UAV height, valid
ground receiver metrics (building pixels excluded from loss and metrics), native
$513 \times 513$ pixel maps, and no test data or field calibration after the
train calibrated priors and the model under evaluation are frozen. Calibration is
therefore not shown as a binary column in Table~\ref{tab:fairness}: the relevant
question is not whether a method has fitted parameters, but whether it adapts
to the evaluation scene using validation, test, or field measurements.
Table~\ref{tab:fairness} summarises the structural axes that can be compared
compactly.

\begin{table}[ht]
\centering
\caption{Evaluation condition comparison across representative methods.}
\label{tab:fairness}
\small
\begin{tabularx}{\textwidth}{Xcccc}
\toprule
\textbf{Method} & \textbf{Unseen city} & \textbf{Cont.\ UAV height}
  & \textbf{$\ge$512\,px} & \textbf{Comparable?} \\
\midrule
PMNet / ICASSP~2023 \cite{pmnet_icassp2023} & No & No & No & No \\
RMTransformer \cite{rmtransformer2025} & No & No & No & No \\
RadioGUNet \cite{radiogunet2025} & Partial & No & No & Partial \\
AIRMap \cite{airmap2025} & Partial & No & No & Partial \\
Gao et al.\ \cite{gao2026} & No & No & Partial & Partial \\
Tarhouni et al.\ \cite{tarhouni2025} & No & No & No & Physical ref. \\
PathFinder \cite{pathfinder2025} & Partial & No & No & Partial \\
This thesis task contract & Yes & Yes & Yes & Reference contract \\
\bottomrule
\end{tabularx}
\end{table}

Two consequences follow from this checklist. First, method rankings can change
substantially when the split changes from random maps to unseen cities. A model
that has learned the texture, street orientation, and height statistics of a
known city may still fail when the same transmitter height range is evaluated
on a new morphology. For this reason, this thesis treats city holdout as a
core part of the task definition rather than as an optional robustness test.

Second, channel attenuation, delay spread, and angular spread should not be treated as
three copies of the same regression problem. LoS path loss often contains a
strong deterministic radial component, whereas NLoS path loss depends on
blocked link geometry and local morphology. Delay and angular spreads are even
more distributional: many pixels are near zero, but physically important tails
remain. A fair comparison therefore has to report the target, regime, masking
policy, and split together; otherwise an apparently lower RMSE can simply mean
that an easier quantity was measured.

\subsubsection{Tarhouni et al.: measured suburban path loss prediction}

Tarhouni et al.\ \cite{tarhouni2025} study machine learning based path loss
prediction in a sub-\SI{6}{GHz} suburban campus. It is useful because the errors
are in physical dB and tied to a concrete environment, but it is terrestrial,
lower frequency, and not a dense UAV CKM with delay- and angular spread targets.
Its role here is qualitative: real environment generalisation tends to move
back toward multi dB errors, which supports keeping city holdout as part of the
task definition.

\subsubsection{AIRMap: adaptive digital twin radio maps}

AIRMap \cite{airmap2025} is an adaptive U-Net autoencoder processing fixed
\(200\times200\) elevation-map tensors with variable spatial resolution
(2.5--15~m/pixel, corresponding to \SI{500}{m}--\SI{3}{km} per side) for ultra low latency digital twin
synchronisation.  It reports sub-\SI{4}{dB} RMSE against ray traced path gain
maps, and a separate lightweight calibration step using 20\,\% of field
measurements reduces the median measurement domain error to about five percent.
Its strong numbers are not directly comparable to CKM because the target is path
gain rather than absolute path loss, the input is a single elevation channel,
and the transmitter configuration is not the continuous UAV height setting used
in this thesis. Its reported \(\sim\)\SI{4}{ms} inference time on an NVIDIA L40S
is nevertheless a strong deployment reference: the value of a CKM surrogate is
not only accuracy, but producing dense maps fast enough for planning loops.

\subsubsection{Gao et al.: multilayer segmentation with corridor weighting}

Gao et al.\ \cite{gao2026} propose a ResNet based outdoor path loss predictor
with Tx depth maps, Rx depth maps, a distance map, and a weighting map that
emphasises the direct Tx/Rx corridor.  Using the ICASSP~2023 outdoor data,
Table~III of the paper reports their method at \SI{5.59}{dB} RMSE, compared
to \SI{8.09}{dB} for PPNet, \SI{9.56}{dB} for RPNet, and \SI{8.72}{dB} for a
twelve layer Vision Transformer baseline; Table~II reports \SI{12.19}{dB}
RMSE on the ITU Challenge 2024 large-area subset, also ahead of the same
baselines, with roughly 60\% fewer FLOPs. This formalises an internal lesson of
the CKM work: the important geometry is not the whole map uniformly, but the
subset near the propagation corridor.

The lesson for this thesis is methodological: Tx centred distance and
obstruction features are useful but are weaker than a true per receiver
corridor description. A full Gao style per Rx corridor map remains future work
because generating one for every receiver pixel at \(513\times513\) resolution
is substantially more expensive than using Tx centred distance and obstruction
features.

\subsubsection{PathFinder: multi transmitter domain adaptation}

PathFinder \cite{pathfinder2025} tackles distribution shift via disentangled
feature encoding and mask guided low rank attention, reporting the lowest MSE
of~$0.1068$ and RMSE of~$0.3263$ on unseen rural datasets. On the
RadioMap3DSeer DS-RPP split, its normalised RMSE of~$0.0331$ can be
approximately converted with the same \SI{36}{dB} RadioMap3DSeer image window convention,
giving about \SI{1.19}{dB}. The unseen rural score would be about
\SI{11.75}{dB} under the same window, but this is only an approximate scale
reading because the paper and code report normalised PNG image units and do not
state a separate rural physical target window. The explicit
domain generalisation focus aligns with the CKM city holdout spirit, though
PathFinder's setup involves multi transmitter inputs and a different
normalisation convention. Techniques from PathFinder such as
transmitter oriented mixup augmentation are listed as medium priority future
work. The same concern motivates the city holdout split used in this thesis: a
CKM surrogate is useful only if it works on new urban layouts, not only on
shuffled samples from the same city.

\subsubsection{ICASSP 2025 indoor challenge}

The ICASSP 2025 Indoor Pathloss Challenge \cite{icassp2025indoor} presents a
structurally different problem, dense wall penetration NLoS effects in bounded
interior spaces, but its public results are a useful warning against
overinterpreting sub-\SI{2}{dB} outdoor challenge numbers.  The official result
page reports that the winning SIP2Net submission obtains a weighted average RMSE
of \SI{9.411}{dB}; IPP-Net reaches \SI{9.501}{dB}; TerRaIn reaches
\SI{10.325}{dB}; and TransPathNet reaches \SI{10.397}{dB}
\cite{indoor2025results,sip2net2025,ippnet2025,transPathNet2025}.  These are
not outdoor UAV results, but they show that when NLoS physics becomes harder
and less directly observable from the input map, modern neural architectures
again hit a roughly \SIrange{9}{10}{dB} error band. For this thesis, the
lesson is diagnostic rather than competitive: harder NLoS settings favour
strong priors, regime separation, and residual distribution modelling over
plain direct regression.

% Section separator
\subsection{Emerging techniques and implications for this thesis}
\label{sec:soa_emerging}
\label{sec:soa_position}
% Section separator

Across the benchmarks reviewed here, the recurring $9$ to
$10\,\mathrm{dB}$ error band in harder NLoS heavy settings reflects the
difficulty of inferring volumetric propagation from a 2D building height map.
The most plausible routes beyond that wall are richer geometry and
diffraction priors \cite{geomDL2024}, PDE style physics losses such as ReVeal
\cite{reveal2025}, sparse target domain adaptation such as RadioPiT
\cite{radiopit2025}, diffusion refinement \cite{radiodiff2025}, and larger
radio map foundation models \cite{fmrme2026}.

For the present thesis, the literature points to four practical design
requirements. First, the comparison contract must be explicit because datasets,
splits, height protocols, and target units differ widely. Second, A2G prediction
requires continuous UAV height conditioning rather than a fixed transmitter
height. Third, LoS/NLoS and environment type effects should be reported
separately because one global RMSE hides the hardest regimes. Fourth,
distribution aware outputs are useful when targets contain multiple modes or
long tails. These observations motivate the methodology in Chapter~3.
